%=========================================================================
\chapter{Basic Set and Map Data Structures}
\label{chap:setsmaps}
\renewcommand{\sectiontitle}{}
%=========================================================================

The \glspl{API} of the various classes presented in the previous \crefrange{chap:swmodel}{chap:sysmodel} regularly use sets containing \gls{WIR} objects like, e.g., \texttt{WIR\_BasicBlockSet}, \texttt{WIR\_SectionSet} or \texttt{WIR\_SymbolSet}.
In general, sets of objects as well as maps that associate specific additional data with objects used as keys play an important role.
As the \gls{WIR} library is implemented in C++17, it is natural to use standard \gls{STL} containers like \texttt{std::set} and \texttt{std::map} for this purpose.
Of course, such containers should not contain the actual \gls{WIR} objects as this would result in severe copying overheads and the risk of corrupt data structures.
Thus, a straightforward realization of, e.g., a set of basic blocks would be to use a set of pointers to basic blocks such as \texttt{std::set<WIR\_BasicBlock *>}.
However, \gls{WIR} is designed to never publicly expose such kinds of data structures for the following reasons:
\begin{itemize}
  \item As motivated in \cref{sec:keyfeatures} on \cpageref{sec:keyfeatures}, the \gls{API} of \gls{WIR} is intentionally kept free of pointers as much as possible.
  Instead, wrapped references are heavily used in order to avoid memory mismanagement, which is why the standard class \texttt{std::reference\_wrapper} regularly appeared in earlier chapters.
  \item Such sets or maps of \gls{WIR} objects are regularly traversed and iterated.
  This, however, is non-deterministic if pointers are used as keys.
  Internally, these \gls{STL} sets and maps are ordered by the actual addresses of the pointers.
  These addresses are determined by the host system's operating system and cannot be controlled within the \gls{WIR} library.
  Thus, depending on whether a \gls{WIR} application gets executed on this or the other host machine, or whether the host's current memory usage is high or low, completely different addresses are assigned to \gls{WIR} objects.
  This may lead to varying, non-deterministic orders of objects within their sets or maps and thus in different orders how \gls{WIR} iterates and thus processes them.
  Such behavior opens the door to numerous kinds of undesired behaviors and bugs that are extremely hard to debug.
\end{itemize}

In order to address the issues above, \gls{WIR} features deterministic mechanisms that allow to obtain a unique numerical identifier for (almost) all kinds of \gls{WIR} objects.
These identifiers and their underlying \gls{API} are specified in \cref{sec:ids}.
\cref{sec:determinism} describes the way how these IDs are used to declare deterministic set- and map-based data structures.


%-------------------------------------------------------------------------
\section{Unique Numerical Identifiers}
\label{sec:ids}
%-------------------------------------------------------------------------

In a straightforward data structure like, e.g., the abovementioned \texttt{std::set<WIR\_BasicBlock *>}, the actual address of a pointer serves as unique key.
This bases on the assumption that \gls{WIR} objects such as basic blocks all have different and unique addresses in memory so that they are uniquely identified by their pointers' addresses.
In order to remove pointers as unique identifiers from deterministic \gls{WIR} data structures, some other mechanism is required that allows to clearly distinguish between two \gls{WIR} objects and to arrange them properly in ordered data structures such as sets or maps.

This mechanism providing unique numerical identifiers for \gls{WIR} objects is realized by class \texttt{WIR\_ID\_API}.
This is a very simple base class from which all the classes presented in \crefrange{chap:lists}{chap:sysmodel} (and many others) publicly inherit so that its \gls{API} shown in \cref{lst:wiridapi} is available (almost) everywhere within \gls{WIR}.
This class internally maintains a static counter that is incremented every time that a new object of a derived class is constructed.
The current counter value is stored for each newly constructed object, thus representing a unique ID of each \gls{WIR} object.

The data type \texttt{WIR\_id\_t} that is used everywhere within \gls{WIR} to represent IDs is declared to be \texttt{unsigned long long} (cf. \cref{lst:wiridapi1}).
Two operators \texttt{==} and \texttt{!=} for equality and inequality tests are declared in \cref{lst:wiridapi2,lst:wiridapi3}.
These operators do not do any kind of semantical test whether two \gls{WIR} objects are equal or inequal.
Instead, only the IDs of both objects are compared.
If this behavior is undesired, both operators can be overridden in derived classes, since they are declared to be virtual.
Finally, method \texttt{getID} (\cref{lst:wiridapi4}) can be used to retrieve an object's unique identifier.

\begin{lstlisting}[style=code,float=t,caption={Handling of Unique Object Identifiers by Class \texttt{WIR\_ID\_API}},label=lst:wiridapi]
// Type WIR_id_t represents unique numerical IDs of WIR objects.
using WIR_id_t æ=æ unsigned long long;Ü\label{lst:wiridapi1}Ü

// Determine whether the IDs of two WIR objects are equal.
virtual @bool@ WIR_ID_API::operator æ==æ @(@ const WIR_ID_API æ&æi @)@ const;Ü\label{lst:wiridapi2}Ü

// Determine whether the IDs of two WIR objects are inequal.
virtual @bool@ WIR_ID_API::operator æ!=æ @(@ const WIR_ID_API æ&æi @)@ const;Ü\label{lst:wiridapi3}Ü

// Get a WIR object's unique ID.
WIR_id_t WIR_ID_API::getID@( void )@ const;Ü\label{lst:wiridapi4}Ü
\end{lstlisting}


%-------------------------------------------------------------------------
\section{Deterministic WIR Data Structures}
\label{sec:determinism}
%-------------------------------------------------------------------------

Whenever sets of \gls{WIR} objects have been introduced in earlier chapters, this guide on purpose used the formulation ``is basically defined as \texttt{std::set<std::reference\_wrapper<...>{}>} within \gls{WIR}'', see e.g., the running example \texttt{WIR\_BasicBlockSet} on \cpageref{def:bbset}.
This, however, is not the complete truth, as this formulation lacks any implementation details making such a set (or map) deterministic.

\begin{lstlisting}[style=code,float=!b,caption={Declaration of Deterministic Data Structures in file \texttt{wirtypes.h}},label=lst:wirsets]
// Comparison class to sort sets and maps of WIR objects uniquely by their IDs.
templateæ<ætypename Tæ>æ struct WIR_CompareÜ\label{lst:wirsets1}Ü
æ{æ
  inline @bool@ operator@()(@ const ästd::reference_wrapperäæ<æTæ> &ælhs, const ästd::reference_wrapperäæ<æTæ> &ærhs @)@ constÜ\label{lst:wirsets2}Ü
  æ{æ
    return@(@ lhs.get@()@.getID@()@ < rhs.get@()@.getID@()@ @)@;Ü\label{lst:wirsets3}Ü
  æ}æ;
æ}æ;Ü\label{lst:wirsets4}Ü

// WIR_BasicBlockSet represents deterministic sets of references to WIR basic blocks.
using WIR_BasicBlockSet æ=æ ästd::setäæ<æästd::reference_wrapperäæ<æWIR_BasicBlockæ>æ, WIR_Compareæ<æWIR_BasicBlockæ>>æ;Ü\label{lst:wirsets5}Ü

// WIR_SectionSet represents deterministic sets of references to WIR sections.
using WIR_SectionSet æ=æ ästd::setäæ<æästd::reference_wrapperäæ<æWIR_Sectionæ>æ, WIR_Compareæ<æWIR_Sectionæ>>æ;Ü\label{lst:wirsets6}Ü

// WIR_SymbolSet represents deterministic sets of references to WIR symbols.
using WIR_SymbolSet æ=æ ästd::setäæ<æästd::reference_wrapperäæ<æWIR_Symbolæ>æ, WIR_Compareæ<æWIR_Symbolæ>>æ;Ü\label{lst:wirsets7}Ü

...
\end{lstlisting}

To achieve this property, \gls{WIR} exploits that the standard \gls{STL} containers for sets and maps can be used in combination with a custom C++ comparison class that in turn determines the ordering of elements within a set or map.
This comparison class as well as all basic deterministic data structures of \gls{WIR} are declared centrally in the source directory in file \texttt{WIR/wir/wirtypes.h}.
An excerpt of this header file is given in \cref{lst:wirsets}.

The comparison class shown in \crefrange{lst:wirsets1}{lst:wirsets4} is one of the very few cases of a template-based class within \gls{WIR}.
It allows to to perform generic less-than comparisons for any kind of class \texttt{T} that is derived from \texttt{WIR\_ID\_API}.
This comparison class gets two (wrapped) references as argument in \cref{lst:wirsets2} and simply performs a numerical less-than comparison of the objects' unique and deterministic IDs (\cref{lst:wirsets3}).

This comparison class is then used in the declarations of \gls{WIR} sets of, e.g., basic blocks, sections or symbols.
This is shown in \crefrange{lst:wirsets5}{lst:wirsets7} where the generic comparison class is specialized towards, e.g., class \texttt{WIR\_BasicBlock}.
The \gls{STL} containers resulting from these declarations now internally sort all their entries according to their IDs which, as motivated in \cref{sec:ids}, no longer on \gls{WIR}-external effects such as host memory addresses.

The entire \gls{WIR} library internally uses only such deterministic data structures whenever sets or maps need to be iterated.
It is good practice that application code using \gls{WIR} also follows this scheme and uses the comparator class from \cref{lst:wirsets}.


\cleardoubleoddpage
